\documentclass[journal, a4paper]{IEEEtran}
\usepackage{booktabs}
\usepackage{graphicx}
\usepackage{stfloats}
\usepackage{placeins}
\usepackage{url}
\usepackage{cite}
\usepackage{cuted}
\usepackage{amssymb}
\usepackage{lipsum}
\usepackage{amsmath}
\usepackage{siunitx}
\sisetup{output-exponent-marker=\ensuremath{\mathrm{e}}}

\begin{document}

\title{Effective Models of Traversable Wormhole Mouths: Fields, Momentum Transfer, Polarization, and Time Slip}

\author{Jan~Ol\v{s}ina\\[1ex] Independent Researcher, Prague, Czech Republic\\
email: \texttt{jan.olsina@gmail.com},  ORCID: \texttt{0009-0003-9816-3958}.}

\maketitle

\begin{abstract}
This note collects and systematizes a sequence of exploratory calculations concerning an idealized traversable wormhole with two compact mouths. The purpose is not to claim a new exact solution of the Einstein equations, but to build a transparent hierarchy of models in which questions about fields, induced monopoles, momentum transfer, mouth orientation, and time offsets can be stated precisely. The simplest model excises two spherical regions and identifies their boundaries through a zero-length throat. In this setting, electrostatic and Newtonian scalar fields are described by boundary matching conditions, while matter crossing the throat is mapped by a local Lorentz transformation between mouth frames. Exact image constructions are reviewed for a conducting sphere and for a simple spherical wormhole; the latter generally requires a distributed image rather than a single Kelvin point charge. Cubical mouths are discussed as a piecewise-flat variant in which crossing kinematics simplifies but the Green function becomes a boundary-integral problem with edge and corner singularities. A weak-field gravitational polarization model is then derived, including the equal-and-opposite apparent mouth monopoles induced by unequal external potentials. The relation of this model to general relativity is assessed using weak-field expansions and junction conditions. Finally, a constant time slip between static mouths is represented as temporal holonomy. A smooth finite-length Ellis-type throat is used to show explicitly that a constant time slip can be introduced without additional local curvature or test-particle force, even though it can have important global, causal, wave-interference, and self-interaction consequences. Throughout, results that follow from geometry or conservation laws are distinguished from assumptions about the exotic material supporting the throat.
\end{abstract}

\begin{IEEEkeywords}
wormholes, traversable wormholes, image charges, self-force, Newtonian limit, thin shells, temporal holonomy, time machines, effective field theory.
\end{IEEEkeywords}

\section{Purpose and Status of the Note}

This document is a structured record of an exploratory model-building discussion. It is intended primarily as a note to the author and to other interested readers. It is not presented as a complete theory of traversable wormholes, nor as a proof that the assumed mouths can be constructed from physically reasonable matter.

The guiding strategy is to separate four logically distinct ingredients:
\begin{enumerate}
    \item the topology and geometry used to identify the two mouths;
    \item the local equations obeyed by matter and fields away from the throat;
    \item matching conditions imposed at the throat;
    \item the constitutive law of the material or field that supports the throat.
\end{enumerate}
The first three ingredients determine many kinematic and conservation statements. The fourth is needed to determine stability, recoil, deformation, energy storage, and the detailed response of the throat.

Traversable wormholes in classical general relativity usually require stress--energy violating the standard pointwise energy conditions \cite{MorrisThorne1988,Visser1995}. Thin-shell constructions make this requirement especially explicit through the junction conditions \cite{Israel1966,PoissonToolkit}. The effective models below should therefore be read as boundary theories whose microscopic support has intentionally been left unspecified.

\section{Geometric Model of Two Mouths}

\subsection{Spherical zero-length identification}

Let two spherical mouths of radius \(a\) have centers \(\mathbf X_A(t)\) and \(\mathbf X_B(t)\), with body-frame orientations \(R_A(t)\) and \(R_B(t)\). Remove the interiors of the two spheres from an otherwise weakly curved external space. A point on mouth \(A\) can be written
\begin{equation}
    \mathbf x_A=\mathbf X_A+aR_A\mathbf n,
    \qquad
    |\mathbf n|=1.
\end{equation}
The corresponding point on mouth \(B\) is
\begin{equation}
    \mathbf x_B=\mathbf X_B+aR_BP\mathbf n,
\end{equation}
where \(P\) is an orthogonal map describing the internal orientation of the throat. The relative orientation seen in the external frame is
\begin{equation}
    \mathcal O=R_BPR_A^{-1}.
\end{equation}

A spacetime version of the gluing rule is
\begin{equation}
    (\tau,\Omega)_A\sim(\tau,\mathcal O\Omega)_B,
    \label{eq:basic_identification}
\end{equation}
where \(\Omega=(\theta,\varphi)\) labels a point on the mouth and \(\tau\) is an internal throat-time label.

Equation~\eqref{eq:basic_identification} is an \emph{equivalence relation}, not a homotopy. Before taking the quotient, the two labels denote different boundary events. After taking the quotient, they denote one event of the completed spacetime. A worldline that looks discontinuous in a single exterior coordinate chart can therefore be continuous in the quotient spacetime.

\subsection{Finite throat}

A smooth finite throat may instead be described by a proper longitudinal coordinate \(l\):
\begin{equation}
    \mathrm{d}s^2=
    -N^2(l)c^2\mathrm{d}\tau^2+
    \mathrm{d}l^2+
    r^2(l)\mathrm{d}\Omega^2.
    \label{eq:generic_throat}
\end{equation}
The function \(r(l)\) has a positive minimum at the throat. The two mouths may be attached at \(l=-\ell/2\) and \(l=+\ell/2\). The zero-length model is an idealized limiting boundary theory, not the only possible geometry.

\section{Kinematics of Traversal}

\subsection{Four-momentum map}

Let \(e_A{}^{\hat a}{}_\mu\) and \(e_B{}^\mu{}_{\hat a}\) be orthonormal tetrads carried by the mouths. Let \(W^{\hat a}{}_{\hat b}\) be the Lorentz transformation implemented by the internal throat map. Then an incoming four-momentum transforms as
\begin{equation}
    p_{\mathrm{out}}^\mu=
    e_B{}^\mu{}_{\hat a}
    W^{\hat a}{}_{\hat b}
    e_A{}^{\hat b}{}_\nu
    p_{\mathrm{in}}^\nu.
    \label{eq:four_momentum_map}
\end{equation}
Because \(W\) is Lorentz,
\begin{equation}
    p_{\mathrm{out}}^\mu p^{\mathrm{out}}_\mu
    =
    p_{\mathrm{in}}^\mu p^{\mathrm{in}}_\mu.
\end{equation}
The rest mass of a freely traversing object is therefore unchanged in the ideal elastic model.

For static synchronized mouths,
\begin{equation}
    W=
    \begin{pmatrix}
    1&0\\
    0&\mathcal O
    \end{pmatrix},
\end{equation}
and
\begin{equation}
    E_{\mathrm{out}}=E_{\mathrm{in}},
    \qquad
    \mathbf p_{\mathrm{out}}=\mathcal O\mathbf p_{\mathrm{in}}.
\end{equation}

\subsection{Impulse delivered to the wormhole}

In a common approximately Minkowskian external frame, conservation of total four-momentum gives
\begin{equation}
    \Delta P^\mu_{\mathrm{wh}}
    =
    p^\mu_{\mathrm{in}}-p^\mu_{\mathrm{out}}.
    \label{eq:wormhole_impulse}
\end{equation}
A useful localized bookkeeping convention is
\begin{equation}
    \Delta P_A^\mu\simeq p_{\mathrm{in}}^\mu,
    \qquad
    \Delta P_B^\mu\simeq-p_{\mathrm{out}}^\mu.
\end{equation}
Only the sum in \eqref{eq:wormhole_impulse} is model-independent. The exact division between mouth recoil, throat stress, radiation, and external supports depends on the throat action.

For a static spatial rotation through an angle \(\alpha\),
\begin{equation}
    \left|\Delta\mathbf P_{\mathrm{wh}}\right|
    =
    2p\sin\left(\frac{\alpha}{2}\right).
\end{equation}
A reversal gives \(2p\); identical input and output directions give zero net impulse to the pair, although the two mouths still receive opposite local impulses.

\subsection{Moving and rotating mouths}

In the nonrelativistic limit, let the surface velocities at the crossing points be
\begin{equation}
    \mathbf v_i^{\mathrm{surf}}
    =
    \mathbf V_i+\boldsymbol\Omega_i\times\mathbf r_i.
\end{equation}
The natural elastic map is
\begin{equation}
\begin{split}
    \mathbf u_{\mathrm{out}}
    =
    \mathbf v_B^{\mathrm{surf}}
    +\mathcal O\left(
    \mathbf u_{\mathrm{in}}-\mathbf v_A^{\mathrm{surf}}
    \right).
    \label{eq:moving_mouth_velocity}
\end{split}
\end{equation}
The mechanical energy transferred to the wormhole and its actuators is
\begin{equation}
    \Delta E_{\mathrm{wh}}
    =
    \frac{m}{2}
    \left(
    |\mathbf u_{\mathrm{in}}|^2-
    |\mathbf u_{\mathrm{out}}|^2
    \right).
\end{equation}
Thus moving mouths can accelerate a particle, but the energy comes from mouth kinetic energy, an actuator, or an internal throat reservoir.

\section{Classical Effective Action}

A broad relativistic action is
\begin{equation}
\begin{split}
S={}&
\frac{c^3}{16\pi G}
\int R\sqrt{-g}\,\mathrm{d}^4x
-\frac{1}{4\mu_0}
\int F_{\mu\nu}F^{\mu\nu}\sqrt{-g}\,\mathrm{d}^4x\\
&+S_{\mathrm{support}}
-mc\int\mathrm{d}s
+q\int A_\mu\mathrm{d}x^\mu.
\label{eq:relativistic_action}
\end{split}
\end{equation}
The unknown term \(S_{\mathrm{support}}\) holds the throat open and determines its mechanical response.

A quasistatic weak-field analog uses an electrostatic potential \(\phi\) and a Newtonian potential \(\Phi\):
\begin{equation}
\begin{split}
L={}&
\frac{m}{2}\dot{\mathbf x}^{\,2}
+\sum_{i=A,B}
\left[
\frac{M_i}{2}\dot{\mathbf X}_i^{\,2}
+\frac{1}{2}\boldsymbol\Omega_i^TI_i\boldsymbol\Omega_i
\right]
-U_{\mathrm{mount}}\\
&+\int_\Omega
\left[
\frac{\epsilon_0}{2}|\boldsymbol\nabla\phi|^2
-\frac{1}{8\pi G}|\boldsymbol\nabla\Phi|^2
\right]\mathrm{d}^3x
-q\phi(\mathbf x)-m\Phi(\mathbf x).
\label{eq:weak_action}
\end{split}
\end{equation}
Variation gives
\begin{equation}
    \nabla^2\phi=-\frac{\rho_e}{\epsilon_0},
    \qquad
    \nabla^2\Phi=4\pi G\rho_m.
\end{equation}
The unusual information is contained in the topology and boundary matching, not in these local equations.

\section{Potentials, Flux, and Self-Interaction}

\subsection{Potential matching}

For a synchronized static throat whose induced metric and gauge potential are continuous, corresponding mouth points satisfy
\begin{equation}
    \phi_A(\Omega)=\phi_B(\mathcal O\Omega).
    \label{eq:electrostatic_matching}
\end{equation}
In the weak-field metric
\begin{equation}
    g_{tt}\simeq-\left(1+\frac{2\Phi}{c^2}\right),
\end{equation}
continuity of the lapse similarly gives
\begin{equation}
    \Phi_A(\Omega)=\Phi_B(\mathcal O\Omega)
\end{equation}
at leading order.

This statement requires a globally exact static potential and excludes a distributional voltage or redshift layer at the throat. Multiply connected spaces can also support source-free harmonic modes. In electrostatics these modes are labeled by the flux threading the throat, an example of Wheeler's ``charge without charge'' idea. Krasnikov separates the charge-generated and source-free components explicitly \cite{Krasnikov2008}.

\subsection{Equal potential does not imply zero force}

The field Green function may be decomposed as
\begin{equation}
    G_{\mathrm{wh}}(\mathbf x,\mathbf x')
    =
    G_0(\mathbf x,\mathbf x')
    +G_{\mathrm{top}}(\mathbf x,\mathbf x').
\end{equation}
After subtraction of the local Coulomb singularity, a charge has regularized self-energy
\begin{equation}
    U_{\mathrm{self}}(\mathbf x)
    =
    \frac{q^2}{2}G_{\mathrm{top}}(\mathbf x,\mathbf x),
\end{equation}
and self-force
\begin{equation}
    \mathbf F_{\mathrm{self}}
    =
    -\boldsymbol\nabla U_{\mathrm{self}}.
\end{equation}
Explicit wormhole calculations find nonzero attractive electromagnetic self-forces even in smooth or short-throat geometries \cite{KhusnutdinovBakhmatov2007,Krasnikov2008}.

\section{Conducting-Sphere Analog}

\subsection{Grounded sphere}

Let a charge \(q\) be at
\begin{equation}
    \mathbf x_0=D\hat{\mathbf n},
    \qquad D>a,
\end{equation}
outside a grounded conducting sphere of radius \(a\). The exact Kelvin image is
\begin{equation}
    \mathbf x'=\frac{a^2}{D}\hat{\mathbf n},
    \qquad
    q'=-q\frac{a}{D}.
    \label{eq:kelvin_image}
\end{equation}
The exterior potential is
\begin{equation}
    \phi(\mathbf x)
    =
    \frac{1}{4\pi\epsilon_0}
    \left[
    \frac{q}{|\mathbf x-\mathbf x_0|}
    +
    \frac{q'}{|\mathbf x-\mathbf x'|}
    \right].
\end{equation}
The physical induced surface density is
\begin{equation}
    \sigma(\theta)
    =
    -\frac{q}{4\pi a}
    \frac{D^2-a^2}
    {\left(D^2+a^2-2aD\cos\theta\right)^{3/2}}.
\end{equation}
For an isolated neutral sphere, a central image \(+qa/D\) is added.

\subsection{Two connected conducting spheres}

For two equal conducting spheres, charges and potentials are related by a capacitance matrix,
\begin{equation}
    \begin{pmatrix}
    Q_A\\Q_B
    \end{pmatrix}
    =
    \begin{pmatrix}
    C_{11}&C_{12}\\
    C_{12}&C_{11}
    \end{pmatrix}
    \begin{pmatrix}
    V_A\\V_B
    \end{pmatrix}.
\end{equation}
An electrical connection imposes \(V_A=V_B\). For \(L\gg a\), the monopole approximation gives
\begin{equation}
    V_A\simeq
    \frac{1}{4\pi\epsilon_0}
    \left(
    \frac{Q_A}{a}+\frac{Q_B}{L}
    \right),
\end{equation}
with an analogous expression for \(V_B\). Equal potentials and equal radii imply \(Q_A=Q_B\).

This conductor analog is useful but stronger than a wormhole boundary condition. A wormhole matches corresponding points; it does not generally make each entire mouth equipotential.

\section{Exact Spherical Wormhole Image Structure}

\subsection{Two flat exteriors glued at a sphere}

Consider two copies of the region \(r>a\), with their \(r=a\) spheres identified. A point charge \(q\) is in the \(+\) exterior at \(r_0>a\). Expand the same-side potential as
\begin{equation}
\begin{split}
4\pi\epsilon_0\phi_+
={}&
\frac{q}{|\mathbf x-\mathbf x_0|}
+\sum_{\ell=0}^{\infty}
\frac{A_\ell}{r^{\ell+1}}
P_\ell(\cos\gamma),
\end{split}
\end{equation}
and the transmitted field as
\begin{equation}
    4\pi\epsilon_0\phi_-
    =
    \sum_{\ell=0}^{\infty}
    \frac{B_\ell}{r^{\ell+1}}
    P_\ell(\cos\gamma).
\end{equation}
Potential and normal-flux matching give
\begin{equation}
    A_\ell
    =
    -\frac{q}{2(\ell+1)}
    \frac{a^{2\ell+1}}{r_0^{\ell+1}},
    \label{eq:spherical_coeff_A}
\end{equation}
\begin{equation}
    B_\ell
    =
    \frac{q(2\ell+1)}{2(\ell+1)}
    \frac{a^{2\ell+1}}{r_0^{\ell+1}},
\end{equation}
for the convention in which both asymptotic potentials vanish and no additional homogeneous flux mode is added.

\subsection{Line-image representation}

Define the Kelvin radius
\begin{equation}
    b=\frac{a^2}{r_0}.
\end{equation}
The coefficients \eqref{eq:spherical_coeff_A} are generated by a uniform fictitious line charge extending from the center to the Kelvin point:
\begin{equation}
    \mathbf x_{\mathrm{im}}(s)=s\hat{\mathbf n}_0,
    \qquad
    0\leq s\leq b,
\end{equation}
with density
\begin{equation}
    \lambda=-\frac{q}{2a}.
\end{equation}
Indeed,
\begin{equation}
\begin{split}
\int_0^b
\frac{\lambda\,\mathrm{d}s}
{|\mathbf x-s\hat{\mathbf n}_0|}
={}&
-\frac{q}{2a}
\sum_{\ell=0}^{\infty}
\frac{b^{\ell+1}}
{(\ell+1)r^{\ell+1}}
P_\ell(\cos\gamma).
\end{split}
\end{equation}
The total line-image charge is
\begin{equation}
    Q_{\mathrm{line}}=-\frac{qa}{2r_0}.
\end{equation}
A different choice of source-free throat flux may add an opposite central image and asymptotic potential offset. Thus the Kelvin point remains geometrically important, but for the wormhole it is generally the endpoint of a distributed image rather than the location of one point image.

\subsection{Self-energy}

At the source,
\begin{equation}
    \int_0^b\frac{\mathrm{d}s}{r_0-s}
    =
    -\ln\left(1-\frac{a^2}{r_0^2}\right).
\end{equation}
The regularized self-energy depends on the chosen homogeneous monopole convention. In the fixed induced-monopole convention used above, the line contribution is attractive and diverges as the charge approaches the ideal sharp mouth. Smooth-throat calculations remove the distributional mouth while retaining a nonzero topology-sensitive self-force \cite{KhusnutdinovBakhmatov2007}.

\section{Cubical Mouths}

\subsection{Geometry and traversal}

Let a mouth be a cube
\begin{equation}
    C=\{\mathbf y:|y_x|,|y_y|,|y_z|\leq a\}.
\end{equation}
A globally consistent face identification is represented by a signed permutation matrix \(P\). On a face with outward body-frame normal \(\mathbf n_f\), the normal component of a traversing velocity must reverse, while tangential coordinates are mapped by \(P\). The body-frame velocity map is
\begin{equation}
    J_f=P\left(I-2\mathbf n_f\mathbf n_f^T\right).
\end{equation}
For static mouths,
\begin{equation}
    \mathbf p_{\mathrm{out}}
    =
    R_BJ_fR_A^T\mathbf p_{\mathrm{in}}.
\end{equation}
This map is constant across each face, making the kinematics simpler than for a spherical mouth.

The ideal rule is ambiguous exactly on an edge or vertex because multiple face normals are available. A physical model should round the edges or specify an edge-scattering law.

\subsection{Curvature concentration}

Each face is intrinsically flat. The two-dimensional surface curvature is concentrated at the vertices through angle deficits. In the three-dimensional space obtained by gluing two cube exteriors, a transverse section through an edge combines two exterior wedges of angle \(3\pi/2\), producing a cone angle
\begin{equation}
    \alpha_{\mathrm{edge}}=3\pi.
\end{equation}
The resulting angular excess is \(-\pi\). Thus edges as well as vertices are singular in the ideal sharp spatial geometry.

\subsection{Green-function decomposition}

Let \(G_D\) and \(G_N\) be the exterior-cube Dirichlet and Neumann Green functions. For a source on side \(A\), the exact wormhole Green functions satisfy
\begin{equation}
    \boxed{
    G_{AA}=\frac{1}{2}(G_N+G_D),
    \qquad
    G_{BA}=\frac{1}{2}(G_N-G_D).
    }
    \label{eq:DN_decomposition}
\end{equation}
At the boundary, \(G_D=0\) and \(\partial_nG_N=0\), so \eqref{eq:DN_decomposition} enforces potential continuity and opposite normal derivatives.

For an infinite plane,
\begin{equation}
    G_D=G_0-G_{\mathrm{im}},
    \qquad
    G_N=G_0+G_{\mathrm{im}}.
\end{equation}
Therefore,
\begin{equation}
    G_{AA}=G_0,
    \qquad
    G_{BA}=G_{\mathrm{im}}.
\end{equation}
A perfectly flat infinite wormhole face produces no leading same-side image force: the Dirichlet and Neumann image terms cancel. The residual field of a finite cubical mouth is therefore associated with its finite extent, edges, corners, and global topology.

There is no Kelvin inversion preserving a cube. The exact ``image charge'' is a continuous equivalent distribution on the six faces, obtained from boundary integral equations. Near a conducting cube edge, the surface charge behaves generically as an integrable fractional singularity, of the form
\begin{equation}
    \sigma\sim\rho^{-1/3},
\end{equation}
where \(\rho\) is the distance to the edge.

\section{Newtonian Gravitational Polarization}

\subsection{Apparent mouth monopoles}

Let two spherical mouths of radius \(a\) be separated by \(L\gg a\). Let \(U_A\) and \(U_B\) denote the externally imposed Newtonian potentials at their centers, excluding the induced mouth fields.

Introduce apparent monopoles \(m_A\) and \(m_B\):
\begin{equation}
    \delta\Phi_A=-\frac{Gm_A}{r_A},
    \qquad
    \delta\Phi_B=-\frac{Gm_B}{r_B}.
\end{equation}
For a fixed zero source-free throat-flux mode, the newly induced monopoles satisfy
\begin{equation}
    m_A+m_B=0.
    \label{eq:zero_net_monopole}
\end{equation}
Potential continuity gives
\begin{equation}
\begin{split}
U_A-\frac{Gm_A}{a}-\frac{Gm_B}{L}
={}&
U_B-\frac{Gm_B}{a}-\frac{Gm_A}{L}.
\end{split}
\end{equation}
Combining this with \eqref{eq:zero_net_monopole},
\begin{equation}
    \boxed{
    m_A=
    \frac{U_A-U_B}
    {2G(a^{-1}-L^{-1})},
    \qquad
    m_B=-m_A.
    }
    \label{eq:induced_monopoles}
\end{equation}
For \(L\gg a\),
\begin{equation}
    m_A\simeq\frac{a}{2G}(U_A-U_B),
    \qquad
    m_B=-m_A.
\end{equation}

Only the potential difference matters. An additive constant in the Newtonian potential cannot induce the pair.

\subsection{Equal source potentials}

Suppose a point mass \(M\) is at equal distance \(D\) from both mouths. The spherical average of its potential over either mouth equals the potential at the center:
\begin{equation}
    U_A=U_B=-\frac{GM}{D}.
\end{equation}
Consequently,
\begin{equation}
    m_A=m_B=0.
\end{equation}
Higher multipoles may remain if the angular field patterns differ, but there is no induced \(1/r\) monopole.

\subsection{A source near only one mouth}

Let the mass be at distance \(D_A\) from \(A\) and \(D_B\) from \(B\). Then
\begin{equation}
    U_A=-\frac{GM}{D_A},
    \qquad
    U_B=-\frac{GM}{D_B}.
\end{equation}
Equation~\eqref{eq:induced_monopoles} gives
\begin{equation}
    \boxed{
    m_A=
    -\frac{M}{2}
    \frac{D_A^{-1}-D_B^{-1}}
    {a^{-1}-L^{-1}},
    \qquad
    m_B=-m_A.
    }
\end{equation}
If \(D_A=D\), \(D_B\gg D\), and \(L\gg a\),
\begin{equation}
    \boxed{
    m_A\simeq-\frac{Ma}{2D},
    \qquad
    m_B\simeq+\frac{Ma}{2D}.
    }
\end{equation}
The near mouth acquires a negative apparent monopole and the remote mouth an equal positive one. This is a redistribution of gravitational flux in the scalar model, not creation of net mass.

\subsection{Caution about gravitational conductors}

An electrical conductor passively equalizes potential because mobile like charges repel. Positive gravitational mass attracts and clumps. If one writes
\begin{equation}
    m_A=\frac{\mathcal M}{2}+\delta,
    \qquad
    m_B=\frac{\mathcal M}{2}-\delta,
\end{equation}
the Newtonian gravitational self-energy contains
\begin{equation}
    -G\delta^2\left(\frac{1}{a}-\frac{1}{L}\right),
\end{equation}
so equal division is an unstable maximum. A gravitational-conductor response must therefore be imposed by exotic stress or active control.

\section{Coupled Motion Near Two Remote Equal Planets}

Although not central to the later relativistic discussion, the earlier orbital model is recorded here for completeness. Let two fixed planets have equal mass \(M\), and let a light mouth of inertial mass \(m\) orbit each planet. Let the local orbital radii be \(r_1\) and \(r_2\), and assume the planetary systems are mutually very distant.

The external potential mismatch is
\begin{equation}
    \Delta\Psi
    =
    -GM\left(\frac{1}{r_1}-\frac{1}{r_2}\right).
\end{equation}
The constrained induced monopole is
\begin{equation}
    \mu
    =
    -\frac{Ma}{2}
    \left(\frac{1}{r_1}-\frac{1}{r_2}\right).
\end{equation}
A convenient effective polarization term is
\begin{equation}
    V_{\mathrm{pol}}
    =
    \frac{GM^2a}{4}
    \left(
    \frac{1}{r_1}-\frac{1}{r_2}
    \right)^2.
    \label{eq:orbital_polarization}
\end{equation}
This term should be understood as the energy of an actively constrained response, not as the stable minimum energy of passive positive masses.

The effective Lagrangian is
\begin{equation}
\begin{split}
L={}&
\sum_{i=1}^{2}
\left[
\frac{m}{2}|\dot{\mathbf r}_i|^2+
\frac{GMm}{r_i}
\right]\\
&-
\frac{GM^2a}{4}
\left(
\frac{1}{r_1}-\frac{1}{r_2}
\right)^2.
\end{split}
\end{equation}
At leading order each force remains radial, so the two individual orbital angular momenta are conserved. The radial motions are coupled. The invariant symmetric family
\begin{equation}
    r_1(t)=r_2(t)
\end{equation}
has zero induced monopole and consists of ordinary Kepler orbits.

Near equal circular radii \(r_0\), write \(r_i=r_0+x_i\). The differential mode \(x_-=x_1-x_2\) obeys
\begin{equation}
    \ddot x_-+\omega_-^2x_-=0,
\end{equation}
with
\begin{equation}
    \omega_-^2=
    \frac{GM}{r_0^3}
    \left(
    1+\frac{M}{m}\frac{a}{r_0}
    \right).
\end{equation}
The factor \(M/m\) illustrates why a mouth can be dynamically affected by the polarization constraint even if its ordinary gravity is negligible to the planet. This model is not automatically compatible with the equivalence principle because the induced gravitational coefficient was varied while the inertial mouth mass was held fixed.

\section{Relation to General Relativity}

\subsection{Weak-field outer region}

In a stationary weak field,
\begin{equation}
\begin{split}
\mathrm{d}s^2={}&
-\left(1+\frac{2\phi}{c^2}\right)c^2\mathrm{d}t^2\\
&+
\left(1-\frac{2\psi}{c^2}\right)
\delta_{ij}\mathrm{d}x^i\mathrm{d}x^j
+O(c^{-4}).
\end{split}
\end{equation}
For slowly moving matter,
\begin{equation}
    \nabla^2\psi=4\pi G\rho.
\end{equation}
If exterior anisotropic stress is negligible, \(\phi=\psi\), and a slow test body satisfies
\begin{equation}
    \ddot{\mathbf x}=-\boldsymbol\nabla\phi.
\end{equation}
The Newtonian scalar equation is therefore a controlled outer limit of general relativity. It need not be valid uniformly through the exotic throat material.

\subsection{Thin-shell junctions}

In a cut-and-paste wormhole, the induced metric \(h_{ab}\) must be continuous. A jump in extrinsic curvature is supported by a surface stress tensor:
\begin{equation}
    \boxed{
    S_{ab}
    =
    -\frac{c^4}{8\pi G}
    \left(
    [K_{ab}]-h_{ab}[K]
    \right).
    }
    \label{eq:israel}
\end{equation}
This is the Israel junction condition \cite{Israel1966,PoissonToolkit}. It supports:
\begin{itemize}
    \item a distributionally thin throat;
    \item local stress--energy conservation on the shell;
    \item weak exterior fields with local \(1/r\) monopoles.
\end{itemize}
It does not imply a universal gravitational-conductor law.

For static synchronized mouths, equality of the induced lapse gives equal weak-field potentials at corresponding points. However, normal derivatives of the metric are determined by \(S_{ab}\). The flux response, susceptibility, and stability depend on the shell equation of state.

\subsection{GR interpretation of the monopoles}

In a buffer region around each nearly spherical mouth,
\begin{equation}
    \phi_i
    =
    U_i-\frac{G\delta M_i}{r_i}
    +
    \sum_{\ell\geq1}
    \frac{\phi_{i,\ell m}}{r_i^{\ell+1}}
    Y_{\ell m}.
\end{equation}
A stationary spherical vacuum monopole perturbation is a perturbation of the local Schwarzschild mass parameter. Thus the \(1/r\) coefficients have a clear relativistic meaning.

Under the additional assumptions that total asymptotic mass, intrinsic shell mass, radiation content, and the homogeneous throat-flux mode remain fixed, one expects
\begin{equation}
    \delta M_A+\delta M_B=0.
\end{equation}
Exchange symmetry and invariance under adding a constant to \(U\) then imply the linear response
\begin{equation}
    \delta M_A=-\delta M_B
    =
    \chi_0(U_A-U_B).
    \label{eq:GR_susceptibility}
\end{equation}
The ideal scalar model predicts
\begin{equation}
    \chi_0^{\mathrm{ideal}}
    =
    \frac{1}{2G(a^{-1}-L^{-1})}.
\end{equation}
General relativity does not make this coefficient universal. It depends on throat geometry, supporting fields, and their constitutive law.

\subsection{Supported and unsupported claims}

The following statements have a controlled geometric or weak-field basis:
\begin{itemize}
    \item local Poisson behavior outside a weakly gravitating throat;
    \item equality of lapse at synchronized identified points;
    \item local mass multipoles in buffer regions;
    \item total stress--energy conservation;
    \item junction stresses for thin shells.
\end{itemize}

The following are model assumptions rather than generic consequences of the Einstein equations:
\begin{itemize}
    \item a perfectly rigid mouth under arbitrary acceleration or tidal loading;
    \item a universal equal-potential susceptibility;
    \item fixed inertial mass while a gravitational monopole changes;
    \item a unique division of recoil between the mouths;
    \item a stable passive gravitational conductor.
\end{itemize}

\section{Accelerated Mouths and Time Slip}

\subsection{Accumulated proper-time difference}

Let mouth \(A\) remain at rest while \(B\) undergoes an accelerated journey and later returns to rest. If \(\tau\) is the common internal throat label, write
\begin{equation}
    t_A=T_A(\tau),
    \qquad
    t_B=T_B(\tau).
\end{equation}
The external time slip is
\begin{equation}
    \Delta(\tau)=T_B(\tau)-T_A(\tau).
\end{equation}

For a special-relativistic trip with \(A\) at rest and \(\tau\) chosen as the proper time of \(B\),
\begin{equation}
    \Delta
    =
    \int
    \left(\gamma_B-1\right)\mathrm{d}\tau,
\end{equation}
or equivalently
\begin{equation}
    \Delta
    =
    \int
    \left[
    1-\sqrt{1-\frac{v_B^2}{c^2}}
    \right]\mathrm{d}t.
\end{equation}
At low velocity,
\begin{equation}
    \Delta\simeq
    \frac{1}{2c^2}
    \int v_B^2\,\mathrm{d}t.
\end{equation}
In a weak gravitational field,
\begin{equation}
\begin{split}
\tau_A-\tau_B\simeq
\frac{1}{c^2}\int
\left[
\Phi_A-\Phi_B
-\frac{v_A^2-v_B^2}{2}
\right]\mathrm{d}t.
\end{split}
\end{equation}
The standard moving-mouth time-machine construction is discussed by Morris, Thorne, and Yurtsever \cite{MorrisThorneYurtsever1988}.

\subsection{Final static identification}

After the acceleration ends, suppose both mouths are again static with equal clock rates. The gluing relation becomes
\begin{equation}
    \boxed{
    (t,\Omega)_A
    \sim
    (t+\Delta,\mathcal O\Omega)_B,
    }
    \label{eq:time_shift_identification}
\end{equation}
with constant \(\Delta\).

For local static mouth metrics
\begin{equation}
    \mathrm{d}s_i^2
    =
    -N_i^2(r_i)c^2\mathrm{d}t_i^2
    +
    A_i^2(r_i)\mathrm{d}r_i^2+
    r_i^2\mathrm{d}\Omega_i^2,
\end{equation}
pullback through \(t_B=t_A+\Delta\) gives
\begin{equation}
    \mathrm{d}(t_A+\Delta)=\mathrm{d}t_A.
\end{equation}
Therefore the constant shift does not change the induced metric or the local junction equations. Frolov, Krtou\v{s}, and Zelnikov discuss an explicit constant-gap quotient with a locally unchanged metric \cite{Frolov2023}.

\subsection{Temporal holonomy}

On a spatial manifold containing a handle, choose a smooth closed one-form \(\eta\) such that
\begin{equation}
    \mathrm{d}\eta=0,
    \qquad
    \oint_\gamma\eta=1,
\end{equation}
where \(\gamma\) passes through the wormhole and returns through exterior space. A stationary metric carrying the time slip is
\begin{equation}
    \boxed{
    \mathrm{d}s^2=
    -N^2c^2(\mathrm{d}T-\Delta\eta)^2
    +
    h_{ij}\mathrm{d}x^i\mathrm{d}x^j.
    }
    \label{eq:holonomy_metric}
\end{equation}
Locally, \(\eta=\mathrm{d}\chi\), and the transformation
\begin{equation}
    \tau=T-\Delta\chi
\end{equation}
removes the cross term. Globally, no single-valued \(\chi\) exists because
\begin{equation}
    \oint_\gamma\Delta\eta=\Delta.
\end{equation}
The slip is therefore locally pure gauge but globally measurable by comparing the wormhole path with an exterior path.

\section{Smooth Finite-Throat Check}

\subsection{Ellis-type throat}

To test whether the absence of a local force is an artifact of an infinitely short or sharp throat, consider the smooth ultrastatic Ellis geometry \cite{Ellis1973,Bronnikov1973}:
\begin{equation}
    \boxed{
    \mathrm{d}s_0^2
    =
    -c^2\mathrm{d}\tau^2+
    \mathrm{d}l^2+
    (l^2+a^2)\mathrm{d}\Omega^2.
    }
    \label{eq:ellis_metric}
\end{equation}
The throat at \(l=0\) is smooth. The Ricci scalar is
\begin{equation}
    R=-\frac{2a^2}{(l^2+a^2)^2},
\end{equation}
which is finite everywhere.

Choose a smooth interpolation
\begin{equation}
    F(-\infty)=0,
    \qquad
    F(+\infty)=1,
\end{equation}
for example
\begin{equation}
    F(l)=\frac{1}{2}
    \left[
    1+\tanh\left(\frac{l}{L_F}\right)
    \right].
\end{equation}
Define
\begin{equation}
    t=\tau+\Delta F(l).
\end{equation}
Then
\begin{equation}
    \boxed{
    \mathrm{d}s_\Delta^2
    =
    -c^2
    \left[
    \mathrm{d}t-\Delta F'(l)\mathrm{d}l
    \right]^2
    +
    \mathrm{d}l^2+
    (l^2+a^2)\mathrm{d}\Omega^2.
    }
    \label{eq:smooth_slip_metric}
\end{equation}
All components are smooth, and the transition can be spread over an arbitrarily long but finite part of the throat.

\subsection{No new curvature}

Equation~\eqref{eq:smooth_slip_metric} follows from the smooth coordinate transformation
\begin{equation}
    \tau=t-\Delta F(l).
\end{equation}
Therefore all local curvature invariants are unchanged:
\begin{equation}
    R[g_\Delta]=R[g_0],
\end{equation}
\begin{equation}
    R_{\mu\nu}R^{\mu\nu}[g_\Delta]
    =
    R_{\mu\nu}R^{\mu\nu}[g_0],
\end{equation}
and similarly for the Riemann tensor and Einstein tensor.

There is no hidden thin shell, delta-function force, or extra stress localized where \(F'\neq0\). For two disconnected asymptotic regions this transformation is globally a clock-convention change. When the two ends are embedded in the same exterior universe, the invariant information is instead the nontrivial global holonomy of \eqref{eq:holonomy_metric}.

\subsection{Explicit geodesics}

Restrict to the equatorial plane. The geodesic Lagrangian is
\begin{equation}
\begin{split}
2\mathcal L={}&
-c^2
\left(
\dot t-\Delta F'\dot l
\right)^2
+\dot l^2+
(l^2+a^2)\dot\varphi^2.
\end{split}
\end{equation}
The conserved energy and angular momentum are
\begin{equation}
    E=c^2(\dot t-\Delta F'\dot l),
\end{equation}
\begin{equation}
    J=(l^2+a^2)\dot\varphi.
\end{equation}
For a timelike geodesic,
\begin{equation}
    \boxed{
    \dot l^2=
    \frac{E^2}{c^2}
    -c^2
    -\frac{J^2}{l^2+a^2}.
    }
\end{equation}
The radial equation is
\begin{equation}
    \boxed{
    \ddot l=
    \frac{J^2l}{(l^2+a^2)^2}.
    }
\end{equation}
Neither equation depends on \(\Delta\). A purely radial geodesic has \(J=0\) and
\begin{equation}
    \ddot l=0.
\end{equation}
The coordinate time is
\begin{equation}
    t(\lambda)
    =
    \tau(\lambda)+\Delta F[l(\lambda)],
\end{equation}
so a complete traversal gains the coordinate-time offset \(\Delta\) without a momentum kick.

\subsection{Generic smooth static throat}

For the more general metric \eqref{eq:generic_throat}, introduce the same transformation:
\begin{equation}
\begin{split}
\mathrm{d}s_\Delta^2={}&
-N^2(l)c^2
\left[
\mathrm{d}t-\Delta F'(l)\mathrm{d}l
\right]^2\\
&+\mathrm{d}l^2+r^2(l)\mathrm{d}\Omega^2.
\end{split}
\end{equation}
The radial first integral is
\begin{equation}
    \dot l^2=
    \frac{E^2}{N^2c^2}
    -\kappa c^2
    -\frac{J^2}{r^2},
\end{equation}
where \(\kappa=1\) for timelike and \(0\) for null geodesics. The radial acceleration is
\begin{equation}
    \ddot l=
    -\frac{E^2N'}{N^3c^2}
    +
    \frac{J^2r'}{r^3}.
\end{equation}
The ordinary redshift force and angular momentum barrier remain, but they are identical for \(\Delta=0\) and constant \(\Delta\neq0\).

\section{When a Time Slip Can Matter Dynamically}

\subsection{Changing time slip}

During acceleration, the map may be
\begin{equation}
    t_B=f(t_A),
\end{equation}
rather than a constant translation. Then
\begin{equation}
    \mathrm{d}t_B=f'(t_A)\mathrm{d}t_A,
\end{equation}
and induced-metric matching requires a corresponding lapse ratio. A changing time slip,
\begin{equation}
    \dot\Delta\neq0,
\end{equation}
can therefore produce redshift, energy exchange, curvature, and stresses. Once the final map is \(f(t)=t+\Delta\), these rate effects disappear while the accumulated offset remains.

\subsection{Time-dependent external fields}

For a field \(\Psi(t,\mathbf x)\), the boundary condition is
\begin{equation}
    \Psi_A(t,\Omega)
    =
    \Psi_B(t+\Delta,\mathcal O\Omega).
\end{equation}
A Fourier component acquires the phase
\begin{equation}
    \widetilde\Psi_B(\omega)
    =
    e^{i\omega\Delta}
    \widetilde\Psi_A(\omega).
\end{equation}
Static fields have \(\omega=0\) and are insensitive to the constant slip. Radiation, oscillating fields, and interference experiments need not be.

\subsection{Self-interaction}

A structureless test particle neglects its own field and sees no extra local force from constant \(\Delta\). A charged or gravitating body can send a field through the throat and interact with a delayed or advanced copy of that field. Such a force is nonlocal and depends on the global Green function, not on new local curvature.

\subsection{Residual excitations}

The acceleration history can leave:
\begin{itemize}
    \item gravitational radiation;
    \item oscillations of the supporting exotic field;
    \item changes of mouth mass, spin, or shape;
    \item stresses in an external accelerating apparatus.
\end{itemize}
These produce real local forces. They are additional final-state data and are not fixed solely by the accumulated value of \(\Delta\).

\subsection{Causal effects}

If the time shift exceeds the shortest exterior causal travel time between the mouths, a closed causal curve can be formed. This is a global causal effect rather than a local force. The local metric can remain stationary and smooth while the global spacetime fails to admit a global time function \cite{MorrisThorneYurtsever1988,Frolov2023}.

\section{Conclusions}

The main conclusions of this note are as follows.

First, an ideal wormhole can be modeled kinematically by identifying two mouth worldtubes. Matter crossing the throat has its four-momentum transformed by the differential of the identification map. The total impulse received by the wormhole system is the difference between incoming and outgoing four-momentum.

Second, electrostatic and Newtonian analogies are useful but should not be confused with a passive physical material. A spherical conducting boundary has a single Kelvin point image, whereas a simple spherical wormhole generally has a distributed line image ending at the Kelvin point. A cubical mouth has simple piecewise-linear traversal kinematics but no finite Kelvin image system; its fields are naturally described by Dirichlet--Neumann Green-function combinations and boundary integral densities.

Third, unequal external scalar potentials can induce equal-and-opposite apparent mouth monopoles. In the ideal short spherical model,
\begin{equation}
    m_A=-m_B
    =
    \frac{U_A-U_B}
    {2G(a^{-1}-L^{-1})}.
\end{equation}
Equal external potentials produce no monopole. The coefficient is not universal in full general relativity because it depends on the throat-supporting stress and its equation of state.

Fourth, the Newtonian model has a legitimate weak-field interpretation outside the throat, and equal lapse at synchronized identified points follows from metric continuity. The detailed flux response, stability, recoil, and energy accounting require a throat action and cannot be inferred from topology alone.

Finally, a constant time slip between mouths is most naturally represented as a time translation in the gluing map or as temporal holonomy. A smooth finite-length throat can carry this holonomy without additional local curvature or test-particle force. The acceleration episode, a changing slip, residual throat excitations, time-dependent fields, self-interaction, and closed causal curves can nevertheless produce important physical effects. The distinction is between a constant affine translation, whose differential is the identity, and dynamical or Lorentz holonomies that act nontrivially on tangent vectors and energy.

\section*{Acknowledgment}
The author wishes to acknowledge the use of large language models (GPT 5.5) for assistance in improving the language and clarity of the manuscript, and for assisting with code and algebraic checking. The author remains responsible for the claims, derivations, citations, and any remaining errors.

\begin{thebibliography}{99}

\bibitem{MorrisThorne1988}
M. S. Morris and K. S. Thorne,
``Wormholes in spacetime and their use for interstellar travel: A tool for teaching general relativity,''
\emph{American Journal of Physics}, vol. 56, no. 5, pp. 395--412, 1988.

\bibitem{MorrisThorneYurtsever1988}
M. S. Morris, K. S. Thorne, and U. Yurtsever,
``Wormholes, time machines, and the weak energy condition,''
\emph{Physical Review Letters}, vol. 61, no. 13, pp. 1446--1449, 1988.

\bibitem{Visser1995}
M. Visser,
\emph{Lorentzian Wormholes: From Einstein to Hawking}.
Woodbury, NY, USA: American Institute of Physics, 1995.

\bibitem{Israel1966}
W. Israel,
``Singular hypersurfaces and thin shells in general relativity,''
\emph{Il Nuovo Cimento B}, vol. 44, pp. 1--14, 1966;
erratum, vol. 48, p. 463, 1967.

\bibitem{PoissonToolkit}
E. Poisson,
\emph{A Relativist's Toolkit: The Mathematics of Black-Hole Mechanics}.
Cambridge, U.K.: Cambridge University Press, 2004.

\bibitem{Krasnikov2008}
S. Krasnikov,
``Electrostatic interaction of a pointlike charge with a wormhole,''
\emph{Classical and Quantum Gravity}, vol. 25, no. 24, Art. no. 245018, 2008,
arXiv:0802.1358.

\bibitem{KhusnutdinovBakhmatov2007}
N. R. Khusnutdinov and I. V. Bakhmatov,
``Self-force of a point charge in the space-time of a symmetric wormhole,''
\emph{Physical Review D}, vol. 76, Art. no. 124015, 2007,
arXiv:0707.3396.

\bibitem{Frolov2023}
V. P. Frolov, P. Krtou\v{s}, and A. Zelnikov,
``Ring wormholes and time machines,''
\emph{Physical Review D}, vol. 108, Art. no. 024034, 2023,
arXiv:2305.03887.

\bibitem{Ellis1973}
H. G. Ellis,
``Ether flow through a drainhole: A particle model in general relativity,''
\emph{Journal of Mathematical Physics}, vol. 14, no. 1, pp. 104--118, 1973.

\bibitem{Bronnikov1973}
K. A. Bronnikov,
``Scalar-tensor theory and scalar charge,''
\emph{Acta Physica Polonica B}, vol. 4, pp. 251--266, 1973.

\bibitem{Jackson1999}
J. D. Jackson,
\emph{Classical Electrodynamics}, 3rd ed.
New York, NY, USA: Wiley, 1999.

\bibitem{Smythe1989}
W. R. Smythe,
\emph{Static and Dynamic Electricity}, 3rd ed.
New York, NY, USA: Hemisphere, 1989.

\end{thebibliography}

\end{document}
